\chapter{Дешифраторы и шифраторы}
\label{ch:coders}

Мы умеем строить любую комбинационную функцию по таблице и упрощать её на
карте Карно. Но у охранной панели с четырьмя зонами возникает практический
вопрос: как по \emph{номеру} зоны включить именно её лампочку, и как
наоборот, по тому, какая из зон сработала, получить её номер и показать
на цифровом индикаторе? Эти задачи встречаются так часто, что получили
готовые «кирпичи»: дешифратор, шифратор и преобразователи кодов. Их
полезно узнавать с первого взгляда. Начнём с узлов, которые преобразуют один
способ записи числа (код) в другой.

\begin{buildbox}
Вкладка \emph{Обучение}, книга «Выбор и распределение», урок
\emph{«Дешифратор»}: соберите дешифратор $2\to4$ с входом разрешения.
Подайте на вход $a$ по очереди $0,1,2,3$ и запишите в тетрадь, какой бит
выхода $y$ равен единице. Дальше в главе мы выясним, как этот
«кирпич» устроен внутри и где он нужен.
\end{buildbox}

\begin{figure}[H]
\centering
\begin{tikzpicture}[node distance=6mm and 32mm]
  \node[block, minimum width=34mm] (bin) {Двоичный код\\\footnotesize \code{10} (число 2)};
  \node[oblock, minimum width=34mm, right=of bin] (oh) {Унитарный код\\\footnotesize \code{0100} (единица на 2-м месте)};
  \draw[arr] ([yshift=2mm]bin.east) -- node[above, font=\sffamily\small] {дешифратор} ([yshift=2mm]oh.west);
  \draw[arr] ([yshift=-2mm]oh.west) -- node[below, font=\sffamily\small] {шифратор} ([yshift=-2mm]bin.east);
\end{tikzpicture}
\caption{Дешифратор и шифратор выполняют взаимно обратные преобразования}
\end{figure}

\term{Унитарный код} (one-hot, «одна горячая») --- код, в котором из $N$
проводов единица стоит ровно на одном, а номер этого провода и есть
передаваемое значение. Он удобен, чтобы что-то \emph{выбрать}: включить одну из
нескольких ламп, одну из нескольких микросхем памяти.

\section{Дешифратор}

\term{Дешифратор} (decoder) имеет $n$ входов и $2^n$ выходов. Входы задают
двоичный номер, и на выходе с этим номером появляется 1, а на всех остальных
--- 0. На схемах дешифратор обозначают буквами DC.

\begin{figure}[H]
\centering
\begin{minipage}{0.4\textwidth}
\centering
\begin{tikzpicture}
  \draw[thick, fill=FpgaLight] (0,0) rectangle (2.6,3.2);
  \draw[thick] (0.6,0) -- (0.6,3.2); \draw[thick] (2.0,0) -- (2.0,3.2);
  \node[font=\sffamily\bfseries\large] at (1.3,1.6) {DC};
  \foreach \i/\n in {0/0, 1/1} {
    \draw[wire] (-0.6,2.1-\i*0.8) node[left] {$a_\i$} -- (0,2.1-\i*0.8);
    \node[font=\sffamily\small] at (0.3,2.1-\i*0.8) {\pgfmathparse{int(2^\i)}\pgfmathresult};
  }
  \draw[wire] (-0.6,0.4) node[left] {$E$} -- (0,0.4); \node[font=\sffamily\small] at (0.3,0.4) {E};
  \foreach \i in {0,...,3} {
    \draw[wire] (2.6,2.6-\i*0.7) -- (3.2,2.6-\i*0.7) node[right] {$y_\i$};
    \node[font=\sffamily\small] at (2.3,2.6-\i*0.7) {\i};
  }
\end{tikzpicture}
\end{minipage}\hfill
\begin{minipage}{0.56\textwidth}
\centering
\begin{tabular}{ccc|cccc}
\toprule
$E$ & $a_1$ & $a_0$ & $y_0$ & $y_1$ & $y_2$ & $y_3$ \\
\midrule
0 & $\times$ & $\times$ & 0 & 0 & 0 & 0 \\
1 & 0 & 0 & \textbf{1} & 0 & 0 & 0 \\
1 & 0 & 1 & 0 & \textbf{1} & 0 & 0 \\
1 & 1 & 0 & 0 & 0 & \textbf{1} & 0 \\
1 & 1 & 1 & 0 & 0 & 0 & \textbf{1} \\
\bottomrule
\end{tabular}

\smallskip
{\small $\times$ --- «неважно, 0 или 1»}
\end{minipage}
\caption{Дешифратор 2$\to$4 с входом разрешения $E$: обозначение и таблица истинности}
\label{fig:dc24}
\end{figure}

Вход $E$ (Enable, разрешение) позволяет «выключить» дешифратор: при $E=0$
все выходы равны нулю, какой бы номер ни был подан.

Как построить дешифратор из элементов? Каждый выход --- это произведение,
равное единице ровно для одной комбинации входов:
\[
  y_0 = E\,\overline{a_1}\,\overline{a_0}, \quad
  y_1 = E\,\overline{a_1}\,a_0, \quad
  y_2 = E\,a_1\,\overline{a_0}, \quad
  y_3 = E\,a_1\,a_0.
\]

\begin{figure}[H]
\centering
\begin{tikzpicture}
  % входы слева; шины a1, ~a1, a0, ~a0, E идут вниз
  \draw[wire] (-0.4,0) node[left] {$a_1$} -- (1.0,0);
  \node[gNOT] (n1) at (2.2,0) {};
  \draw[wire] (1.0,0) -- (n1.input);
  \draw[wire] (n1.output) -- (3.4,0);
  \draw[wire] (-0.4,-1.0) node[left] {$a_0$} -- (4.6,-1.0);
  \node[gNOT] (n0) at (5.8,-1.0) {};
  \draw[wire] (4.6,-1.0) -- (n0.input);
  \draw[wire] (n0.output) -- (7.0,-1.0);
  \draw[wire] (-0.4,-2.0) node[left] {$E$} -- (8.2,-2.0);
  \foreach \x/\byy/\l in {1.0/0/a_1, 3.4/0/\overline{a_1}, 4.6/-1.0/a_0, 7.0/-1.0/\overline{a_0}, 8.2/-2.0/E} {
    \draw[very thick, FpgaBlue] (\x,\byy) -- (\x,-11.3);
    \node[below, font=\small] at (\x,-11.3) {$\l$};
  }
  \foreach \x/\byy in {1.0/0, 4.6/-1.0, 8.2/-2.0} \fill (\x,\byy) circle (1.3pt);
  % элементы И: y0 = ~a1 ~a0, y1 = ~a1 a0, y2 = a1 ~a0, y3 = a1 a0
  \foreach \i/\xp/\xq in {0/3.4/7.0, 1/3.4/4.6, 2/1.0/7.0, 3/1.0/4.6} {
    \pgfmathsetmacro{\y}{-3.8-\i*2.0}
    \node[gAND, logic gate inputs=nnn, minimum height=13.5mm] (g\i) at (11.0,\y) {}; % fig-ok: 3-input gate, 13.5 mm
    \draw[wire] (g\i.input 1) -- (\xp,0 |- g\i.input 1) node[circle, fill, inner sep=1.1pt] {};
    \draw[wire] (g\i.input 2) -- (\xq,0 |- g\i.input 2) node[circle, fill, inner sep=1.1pt] {};
    \draw[wire] (g\i.input 3) -- (8.2,0 |- g\i.input 3) node[circle, fill, inner sep=1.1pt] {};
    \draw[wire] (g\i.output) -- ++(0.6,0) node[right] {$y_\i$};
  }
\end{tikzpicture}
\caption{Дешифратор 2$\to$4 на элементах И и НЕ}
\label{fig:dc24_gates}
\end{figure}

\begin{figure}[H]
\centering
\begin{tikztimingtable}[timing/slope=0.05, xscale=2.2, yscale=1.4, timing/font=\sffamily\small, timing/rowdist=3.2ex]
  $a_1a_0$ & 2D{00} 2D{01} 2D{10} 2D{11} 2D{00} \\
  $y_0$ & 2H 6L 2H \\
  $y_1$ & 2L 2H 6L \\
  $y_2$ & 4L 2H 4L \\
  $y_3$ & 6L 2H 2L \\
\end{tikztimingtable}
\caption{При переборе номеров единица «переходит» с одного выхода на другой}
\end{figure}

\subsection{Где применяются дешифраторы}

\begin{itemize}
  \item \textbf{Выбор устройства по адресу.} В компьютере к одной шине
        подключено много микросхем памяти и периферии. Дешифратор по
        старшим битам адреса выбирает, какая из них сейчас должна ответить.
  \item \textbf{Управление индикаторами}: включение одного из нескольких
        светодиодов, разрядов индикатора, строк матрицы.
  \item \textbf{Выбор команды в процессоре}: код операции дешифрируется в
        управляющие сигналы «сложить», «прочитать из памяти» и т.\,д.
\end{itemize}

\begin{figure}[H]
\centering
\begin{tikzpicture}[node distance=5mm]
  \node[hblock, minimum height=30mm, minimum width=18mm] (cpu) {Процессор};
  \node[block, right=22mm of cpu, minimum height=26mm] (dc) {DC\\2$\to$4};
  \foreach \i/\t in {0/ПЗУ, 1/ОЗУ, 2/Порт ввода, 3/Порт вывода} {
    \node[oblock, minimum width=24mm, minimum height=7mm] (d\i) at (9.0,1.65-\i*1.1) {\t};
  }
  \draw[arr] (cpu.east) -- node[above, lbl] {A[15:14]} (dc.west);
  \foreach \i in {0,...,3} \draw[arr] ($(dc.north east)!{(\i+0.5)/4}!(dc.south east)$) -- ++(0.6,0) |- (d\i.west);
  \foreach \i in {0,...,3} \node[lbl, above=0mm of d\i.north west, anchor=south west] {CS$_\i$};
\end{tikzpicture}
\caption{Дешифратор адреса: два старших бита адреса выбирают одну из четырёх микросхем}
\end{figure}

\subsection{Наращивание: большой дешифратор из малых}

Вход разрешения $E$ нужен не только чтобы «выключить» дешифратор: с его помощью
из малых дешифраторов собирают большие. Дешифратор $3\to8$ получается из
двух дешифраторов $2\to4$: младшие два разряда номера $a_1a_0$ подаются на
оба, а старший разряд $a_2$ решает, какой из двух «просыпается». Первому
разрешение даёт $\overline{a_2}$, второму --- $a_2$. Тогда при $a_2=0$ работает
первый и зажигает одну из линий $y_0\ldots y_3$, при $a_2=1$ --- второй и
зажигает одну из $y_4\ldots y_7$. Тот же приём (первый уровень выбирает
микросхему, второй --- линию) используется при выборе устройства по адресу.

\begin{figure}[H]
\centering
\begin{tikzpicture}
  \node[block, minimum width=22mm, minimum height=16mm] (d0) at (5,1.3) {DC\\2$\to$4};
  \node[block, minimum width=22mm, minimum height=16mm] (d1) at (5,-1.5) {DC\\2$\to$4};
  \draw[arr] (0,0.1) node[left] {$a_1a_0$} -- (2.4,0.1);
  \fill (2.4,0.1) circle (1.3pt);
  \draw[arr] (2.4,0.1) |- ($(d0.west)+(0,-0.3)$);
  \draw[arr] (2.4,0.1) |- ($(d1.west)+(0,0.3)$);
  \draw[arr] (0,2.9) node[left] {$\overline{a_2}$} -- (5,2.9) -- node[right, lbl] {$E$} (d0.north);
  \draw[arr] (0,-3.3) node[left] {$a_2$} -- (5,-3.3) -- node[right, lbl] {$E$} (d1.south);
  \draw[arr] ($(d0.east)+(0,0)$) -- ++(1.4,0) node[right] {$y_0\ldots y_3$};
  \draw[arr] ($(d1.east)+(0,0)$) -- ++(1.4,0) node[right] {$y_4\ldots y_7$};
\end{tikzpicture}
\caption{Дешифратор $3\to8$ из двух дешифраторов $2\to4$ с входом разрешения}
\label{fig:dc38}
\end{figure}

\begin{diybox}[Дешифратор в Logra]
\textbf{Задание.} \stars{1} Соберите дешифратор $2\to4$ и проверьте все
восемь комбинаций $(E,a)$. \stars{2} Сравните его со своей схемой из
элементов И и НЕ (рисунок~\ref{fig:dc24_gates}): число элементов и
результат. \stars{3} Соберите дешифратор $3\to8$ из двух $2\to4$ по
рисунку~\ref{fig:dc38}; критерий приёмки: при любом из восьми номеров
горит ровно одна линия, при $E=0$ нет ни одной.

\textbf{В Logra.} \stars{1} Урок \emph{«Дешифратор»} (книга «Выбор и
распределение»); пример \code{decoder}. \stars{2} Вкладка \emph{Конструктор}:
блоки И, НЕ. \stars{3} Вкладка \emph{Конструктор}: два блока «Дешифратор» с
включённым разрешением, блок «Срез шины» для выделения $a_1a_0$ и $a_2$
(урок «Шина: срез битов»), элемент НЕ.

\textbf{Ожидаемый результат.} Для номера 5 ($a=101$) загорается $y_5$,
остальные выходы нули; для 2 ($a=010$) загорается $y_2$. Пример
\code{binary-onehot} показывает, что «лишние» номера дают пустую
комбинацию, если их не предусмотреть.

\textbf{Как проверить.} Вкладка \emph{Диаграммы}: задайте $a$ счётчиком
$0$--$7$; единица должна «бежать» по выходам. Вкладка \emph{RTL}: сравните
вашу схему с одним блоком «Дешифратор».
\end{diybox}

\section{Семисегментный индикатор}

Семисегментный индикатор состоит из семи светодиодных полосок, обозначаемых
буквами от $a$ до $g$. Зажигая разные сегменты, можно показать любую цифру.
Узел, который переводит двоичный код цифры в сигналы для сегментов, тоже
называют дешифратором (хотя он выдаёт не унитарный код).

\newcommand{\sevenseg}[8]{%
  \begin{tikzpicture}[scale=0.32, seg/.style={line width=2.6pt, line cap=round}]
    \def\segc##1{\ifnum##1=1 \def\cc{FpgaRed}\else\def\cc{gray!20}\fi}
    \segc{#1}\draw[seg, color=\cc] (0.2,4) -- (1.8,4);   % a
    \segc{#2}\draw[seg, color=\cc] (2,3.8) -- (2,2.2);   % b
    \segc{#3}\draw[seg, color=\cc] (2,1.8) -- (2,0.2);   % c
    \segc{#4}\draw[seg, color=\cc] (0.2,0) -- (1.8,0);   % d
    \segc{#5}\draw[seg, color=\cc] (0,1.8) -- (0,0.2);   % e
    \segc{#6}\draw[seg, color=\cc] (0,3.8) -- (0,2.2);   % f
    \segc{#7}\draw[seg, color=\cc] (0.2,2) -- (1.8,2);   % g
    \node[font=\sffamily\footnotesize] at (1,-1.1) {#8};
  \end{tikzpicture}}

\begin{figure}[H]
\centering
\begin{tikzpicture}[scale=0.8, seg/.style={line width=4pt, line cap=round, FpgaRed}]
  \draw[seg] (0.25,4) -- (1.75,4) node[midway, above=1mm, text=FpgaDark, font=\sffamily\small] {a};
  \draw[seg] (2,3.75) -- (2,2.25) node[midway, right=1mm, text=FpgaDark, font=\sffamily\small] {b};
  \draw[seg] (2,1.75) -- (2,0.25) node[midway, right=1mm, text=FpgaDark, font=\sffamily\small] {c};
  \draw[seg] (0.25,0) -- (1.75,0) node[midway, below=1mm, text=FpgaDark, font=\sffamily\small] {d};
  \draw[seg] (0,1.75) -- (0,0.25) node[midway, left=1mm, text=FpgaDark, font=\sffamily\small] {e};
  \draw[seg] (0,3.75) -- (0,2.25) node[midway, left=1mm, text=FpgaDark, font=\sffamily\small] {f};
  \draw[seg] (0.25,2) -- (1.75,2) node[midway, above=0.5mm, text=FpgaDark, font=\sffamily\small] {g};
\end{tikzpicture}
\hspace{12mm}
\sevenseg1111110{0}\,\sevenseg0110000{1}\,\sevenseg1101101{2}\,\sevenseg1111001{3}\,\sevenseg0110011{4}\,%
\sevenseg1011011{5}\,\sevenseg1011111{6}\,\sevenseg1110000{7}\,\sevenseg1111111{8}\,\sevenseg1111011{9}
\caption{Семисегментный индикатор: обозначение сегментов и изображение цифр}
\end{figure}

\begin{table}[H]
\centering
\caption{Таблица истинности дешифратора для семисегментного индикатора (1 --- сегмент горит)}
\label{tab:7seg}
\small
\begin{tabular}{c|cccc|ccccccc}
\toprule
Цифра & $x_3$ & $x_2$ & $x_1$ & $x_0$ & a & b & c & d & e & f & g \\
\midrule
0 & 0&0&0&0 & 1&1&1&1&1&1&0 \\
1 & 0&0&0&1 & 0&1&1&0&0&0&0 \\
2 & 0&0&1&0 & 1&1&0&1&1&0&1 \\
3 & 0&0&1&1 & 1&1&1&1&0&0&1 \\
4 & 0&1&0&0 & 0&1&1&0&0&1&1 \\
5 & 0&1&0&1 & 1&0&1&1&0&1&1 \\
6 & 0&1&1&0 & 1&0&1&1&1&1&1 \\
7 & 0&1&1&1 & 1&1&1&0&0&0&0 \\
8 & 1&0&0&0 & 1&1&1&1&1&1&1 \\
9 & 1&0&0&1 & 1&1&1&1&0&1&1 \\
\bottomrule
\end{tabular}
\end{table}

Каждый столбец a--g --- это отдельная логическая функция от четырёх
переменных. Её можно упростить по карте Карно и собрать из элементов, как в
главах~\ref{ch:comb} и~\ref{ch:minimize}. Коды $1010$--$1111$ не
являются цифрами, поэтому на них можно ничего не требовать: для
сегментных функций это безразличные состояния, и они делают
минимальные формы заметно короче.

\textbf{Пример: сегмент $e$.} В таблице~\ref{tab:7seg} сегмент $e$ горит на
цифрах $0,2,6,8$. Нарисуем карту для функции $e(x_3,x_2,x_1,x_0)$,
отметив безразличные клетки $10$--$15$ символом $\times$.

\begin{figure}[H]
\centering
\begin{tikzpicture}[cell/.style={draw, minimum size=10mm, font=\sffamily\large}]
  \foreach \c/\lab in {0/00, 1/01, 2/11, 3/10} \node[font=\sffamily\small] at (\c,0.85) {\lab};
  \foreach \r/\lab in {0/00, 1/01, 2/11, 3/10} \node[font=\sffamily\small] at (-0.95,-\r) {\lab};
  \node[font=\sffamily\small, anchor=south east] at (-0.5,0.9) {$x_3x_2\backslash x_1x_0$};  % e: 0,2,6,8 = 1; 1,3,4,5,7,9 = 0; 10..15 = x
  \foreach \r/\c/\v in {0/0/1, 0/1/0, 0/2/0, 0/3/1,
                        1/0/0, 1/1/0, 1/2/0, 1/3/1,
                        2/0/x, 2/1/x, 2/2/x, 2/3/x,
                        3/0/1, 3/1/0, 3/2/x, 3/3/x}
    {\def\cc{white}
     \if 1\v \def\cc{FpgaLight}\fi
     \if x\v \def\cc{gray!25}\fi
     \node[cell, fill=\cc] at (\c,-\r) {\if x\v$\times$\else\v\fi};}
  \foreach \x/\y in {0/0, 3/0, 0/-3, 3/-3}
    \draw[very thick, FpgaRed, rounded corners=4pt] (\x-0.44,\y+0.44) rectangle (\x+0.44,\y-0.44);
  \draw[very thick, FpgaGreen, rounded corners=4pt] (2.62,0.38) rectangle (3.38,-3.38);
  \node[font=\sffamily\small, text=FpgaRed, anchor=west] at (4.4,0.2) {$\overline{x_2}\,\overline{x_0}$ --- четыре угла};
  \node[font=\sffamily\small, text=FpgaGreen, anchor=west] at (4.4,-0.6) {$x_1\overline{x_0}$ --- правый столбец};
  \node[font=\sffamily\small\bfseries, anchor=west] at (4.4,-1.6) {$e=\overline{x_0}\,(\overline{x_2}+x_1)$};
\end{tikzpicture}
\caption{Сегмент $e$ на карте Карно: безразличные клетки позволяют получить две крупные группы по четыре клетки}
\label{fig:seg_e}
\end{figure}

Четыре угла (клетки $0,2,8,10$) дают $\overline{x_2}\,\overline{x_0}$,
правый столбец (клетки $2,6,10,14$) --- $x_1\overline{x_0}$. Складываем и
выносим общий множитель: $e=\overline{x_0}(\overline{x_2}+x_1)$ --- всего три
элемента вместо семи слагаемых по четыре буквы. Проверим на цифрах:
для $4$ ($0100$) получаем $1\cdot(0+0)=0$, и сегмент $e$ не горит; для $6$
($0110$) --- $1\cdot(0+1)=1$, горит. Аналогично сегмент $c$ не горит только
на цифре 2, и его минимальная форма --- $c=x_2+\overline{x_1}+x_0$.

\begin{warnbox}[Активный уровень]
Многие реальные индикаторы имеют общий анод: сегмент горит не при единице
на выходе дешифратора, а при \emph{нуле}. Тогда все семь выходов нужно
инвертировать. Если на плате вместо цифры вы увидели «негатив», вероятная
причина именно в этом, а не в ошибке в таблице.
\end{warnbox}

\begin{diybox}[Цифра на семи сегментах]
\textbf{Задание.} \stars{1} Покажите на индикаторе цифры $0\ldots9$.
\stars{2} Добавьте шестнадцатеричные цифры $A\ldots F$ и посмотрите,
чем их вид отличается от «цифр». \stars{3} Найдите минимальную форму
сегмента $c$ и соберите его из двух элементов: $c=x_2+\overline{x_1}+x_0$.

\textbf{В Logra.} \stars{1} Урок \emph{«Цифра на семи сегментах»} (книга
«Индикаторы»); пример \code{segment-decoder} (библиотека примеров).
\stars{2} Урок \emph{«Цифры от 0 до F»}. \stars{3} Вкладка
\emph{Конструктор}: блок ИЛИ с тремя входами и инвертор.

\textbf{Ожидаемый результат.} Выход семи бит; младший бит --- сегмент
$a$, старший --- $g$. Для цифры 1 получается \code{0000110} (горят $b$ и $c$).
Для кодов $10$--$15$ блок «BCD\,$\to$\,7 сегментов» гасит все
сегменты, то есть выбирает для безразличных состояний нули.

\textbf{Как проверить.} Вкладка \emph{Диаграммы}: перебор $0$--$9$
и сравнение с таблицей~\ref{tab:7seg}. Для самодельного сегмента
$c$ сверьте выход с одним столбцом таблицы на всех десяти цифрах; на
$10$--$15$ результат неважен.
\end{diybox}

\section{Шифратор}

\term{Шифратор} (encoder, на схемах обозначается CD) выполняет обратное
преобразование: из унитарного кода на $2^n$ входах делает $n$-разрядный
двоичный номер той линии, на которой стоит единица.

\begin{figure}[H]
\centering
\begin{minipage}{0.38\textwidth}
\centering
\begin{tikzpicture}
  \draw[thick, fill=FpgaLight] (0,0) rectangle (2.6,3.0);
  \draw[thick] (0.6,0) -- (0.6,3.0); \draw[thick] (2.0,0) -- (2.0,3.0);
  \node[font=\sffamily\bfseries\large] at (1.3,1.5) {CD};
  \foreach \i in {0,...,3} {
    \draw[wire] (-0.6,2.4-\i*0.65) node[left] {$x_\i$} -- (0,2.4-\i*0.65);
    \node[font=\sffamily\small] at (0.3,2.4-\i*0.65) {\i};
  }
  \foreach \i/\w in {0/1, 1/2} {
    \draw[wire] (2.6,1.9-\i*0.9) -- (3.2,1.9-\i*0.9) node[right] {$y_\i$};
    \node[font=\sffamily\small] at (2.3,1.9-\i*0.9) {\w};
  }
\end{tikzpicture}
\end{minipage}\hfill
\begin{minipage}{0.58\textwidth}
\centering
\begin{tabular}{cccc|cc}
\toprule
$x_0$ & $x_1$ & $x_2$ & $x_3$ & $y_1$ & $y_0$ \\
\midrule
1 & 0 & 0 & 0 & 0 & 0 \\
0 & 1 & 0 & 0 & 0 & 1 \\
0 & 0 & 1 & 0 & 1 & 0 \\
0 & 0 & 0 & 1 & 1 & 1 \\
\bottomrule
\end{tabular}
\end{minipage}
\caption{Шифратор 4$\to$2: обозначение и таблица истинности}
\end{figure}

По таблице видно: $y_0$ равен 1, когда активна линия 1 или 3; $y_1$ --- когда
активна линия 2 или 3. Значит, шифратор собирается из одних элементов ИЛИ:
\[
  y_0 = x_1 + x_3, \qquad y_1 = x_2 + x_3.
\]

\begin{figure}[H]
\centering
\begin{tikzpicture}
  \foreach \i in {0,...,3} \draw[very thick, FpgaBlue] (\i*0.7,0.6) node[above, text=FpgaDark] {$x_\i$} -- (\i*0.7,-2.4);
  \node[gOR] (o0) at (4.2,-0.4) {};
  \node[gOR] (o1) at (4.2,-2.0) {};
  \draw[wire] (o0.input 1) -- (0.7,0 |- o0.input 1) node[circle, fill, inner sep=1.1pt] {};
  \draw[wire] (o0.input 2) -- (2.1,0 |- o0.input 2) node[circle, fill, inner sep=1.1pt] {};
  \draw[wire] (o1.input 1) -- (1.4,0 |- o1.input 1) node[circle, fill, inner sep=1.1pt] {};
  \draw[wire] (o1.input 2) -- (2.1,0 |- o1.input 2) node[circle, fill, inner sep=1.1pt] {};
  \draw[wire] (o0.output) -- ++(0.6,0) node[right] {$y_0$};
  \draw[wire] (o1.output) -- ++(0.6,0) node[right] {$y_1$};
  \node[lbl, text=FpgaGray] at (0,-3.1) {не используется};
\end{tikzpicture}
\caption{Шифратор 4$\to$2 из двух элементов ИЛИ}
\end{figure}

\subsection{Приоритетный шифратор}

Что будет, если единица окажется сразу на двух входах, например, если
одновременно нажаты кнопки 1 и 2? Простой шифратор выдаст $y_1y_0 = 11$, то есть
номер 3, --- явную ошибку. Поэтому на практике применяют
\term{приоритетный шифратор}: он выдаёт номер \emph{старшего} из активных
входов, а остальные игнорирует. Дополнительный выход $V$ (valid) показывает,
что хотя бы один вход активен.

\begin{table}[H]
\centering
\caption{Приоритетный шифратор 4$\to$2 ($\times$ --- значение не важно)}
\begin{tabular}{cccc|ccc}
\toprule
$x_3$ & $x_2$ & $x_1$ & $x_0$ & $y_1$ & $y_0$ & $V$ \\
\midrule
0 & 0 & 0 & 0 & 0 & 0 & 0 \\
0 & 0 & 0 & 1 & 0 & 0 & 1 \\
0 & 0 & 1 & $\times$ & 0 & 1 & 1 \\
0 & 1 & $\times$ & $\times$ & 1 & 0 & 1 \\
1 & $\times$ & $\times$ & $\times$ & 1 & 1 & 1 \\
\bottomrule
\end{tabular}
\end{table}

Выражения: $y_1 = x_3 + x_2$, \ $y_0 = x_3 + \overline{x_2}\,x_1$, \
$V = x_3 + x_2 + x_1 + x_0$.

\begin{infobox}[Где применяются шифраторы]
\begin{itemize}
  \item \textbf{Клавиатуры}: номер нажатой клавиши превращается в код.
  \item \textbf{Контроллеры прерываний}: из нескольких одновременных запросов
        процессору передаётся номер самого важного.
  \item \textbf{Датчики положения}: из показаний линейки фотодатчиков
        получают число --- координату.
\end{itemize}
\end{infobox}

Почему именно старший, а не младший? Договорённость произвольна, важно
лишь, чтобы она была \emph{одна и та же} везде: в клавиатуре, в контроллере
прерываний, в охранной панели. В панели это значит: если одновременно
сработали зоны~1 и~3, на индикаторе покажется номер~3, и охранник
пойдёт туда, где «важнее». Выход $V$ нужен, чтобы отличить «сработала
зона~0» от «не сработало ничего»: в обоих случаях $y_1y_0=00$.

\begin{diybox}[Шифратор и приоритет]
\textbf{Задание.} \stars{1} Для входов $x_3x_2x_1x_0=0101$ и $1010$
определите выход приоритетного шифратора ($y_1y_0$, $V$). \stars{2}
Соберите шифратор из элементов ИЛИ по формулам $y_0=x_1+x_3$,
$y_1=x_2+x_3$ и сравните с приоритетным при двух единицах на входе.
\stars{3} Добавьте к приоритетному шифратору выход $V$ из элементов
(формула $V=x_3+x_2+x_1+x_0$).

\textbf{В Logra.} \stars{1} Урок \emph{«Приоритетный шифратор»} (книга
«Выбор и распределение»); пример \code{priority} (библиотека примеров).
\stars{2}, \stars{3} Вкладка \emph{Конструктор}: блоки ИЛИ, НЕ, И.

\textbf{Ожидаемый результат.} Для $0101$: $y_1y_0=10$, $V=1$; для $1010$:
$y_1y_0=11$, $V=1$. Простой шифратор на двух единицах выдаёт
\emph{не тот} номер: например, при $x_1=x_2=1$ получается $11$, то есть
номер~3, которого на входе нет.

\textbf{Как проверить.} Вкладка \emph{Диаграммы}: подайте на $x$ все
16 комбинаций; убедитесь, что при единице в $x_3$ всегда $y=11$,
и что $V=0$ только при $x=0000$.
\end{diybox}

\section{Преобразователи кодов}

Дешифраторы и шифраторы --- частные случаи \term{преобразователей кодов}.
Ещё один полезный пример --- перевод двоичного кода в \term{код Грея}, в
котором соседние числа отличаются ровно одним битом. Это важно для датчиков
угла поворота: при переходе от 3 (\code{011}) к 4 (\code{100}) в двоичном коде
меняются сразу три бита, и в промежуточный момент датчик может выдать
что угодно. В коде Грея меняется только один бит. (Это тот же
механизм, что и «иголки» в главе~\ref{ch:comb}: несколько сигналов меняются
не одновременно, и между ними мелькает ложное значение.)

\begin{figure}[H]
\centering
\begin{minipage}{0.45\textwidth}
\centering
\begin{tabular}{c c c}
\toprule
\textbf{Число} & \textbf{Двоичный} & \textbf{Грея} \\
\midrule
0 & \code{000} & \code{000} \\
1 & \code{001} & \code{001} \\
2 & \code{010} & \code{011} \\
3 & \code{011} & \code{010} \\
4 & \code{100} & \code{110} \\
5 & \code{101} & \code{111} \\
6 & \code{110} & \code{101} \\
7 & \code{111} & \code{100} \\
\bottomrule
\end{tabular}
\end{minipage}\hfill
\begin{minipage}{0.5\textwidth}
\centering
\begin{tikzpicture}
  \draw[very thick, FpgaBlue] (0,0) node[left, text=FpgaDark] {$b_2$} -- (3.4,0) node[right, text=FpgaDark] {$g_2$};
  \node[gXOR] (x1) at (2.4,-1.0) {};
  \node[gXOR] (x0) at (2.4,-2.5) {};
  \fill (0.6,0) circle (1.3pt);
  \draw[wire] (x1.input 1) -- (0.6,0 |- x1.input 1) -- (0.6,0);
  \draw[wire] (x1.input 2) -- (0,0 |- x1.input 2) node[left] {$b_1$};
  \fill (1.0,0 |- x1.input 2) circle (1.3pt);
  \draw[wire] (x0.input 1) -- (1.0,0 |- x0.input 1) -- (1.0,0 |- x1.input 2);
  \draw[wire] (x0.input 2) -- (0,0 |- x0.input 2) node[left] {$b_0$};
  \draw[wire] (x1.output) -- ++(0.6,0) node[right] {$g_1$};
  \draw[wire] (x0.output) -- ++(0.6,0) node[right] {$g_0$};
\end{tikzpicture}

\small $g_i = b_{i+1} \oplus b_i$
\end{minipage}
\caption{Двоичный код и код Грея; преобразователь на элементах «Исключающее ИЛИ»}
\end{figure}

Обратное преобразование идёт «цепочкой»: старший бит остаётся, каждый
следующий равен XOR уже полученного двоичного бита и текущего бита Грея,
$b_{n-1}=g_{n-1}$, $b_i=b_{i+1}\oplus g_i$. Например, код Грея \code{110}
даёт $b_2=1$, $b_1=1\oplus1=0$, $b_0=0\oplus0=0$, то есть двоичное
\code{100} (число~4). Прямое преобразование устроено проще: каждый бит
Грея зависит только от двух соседних двоичных, поэтому все XOR работают
параллельно, без цепочки.

\begin{diybox}[Код Грея]
\textbf{Задание.} \stars{1} Запишите в коде Грея числа $0\ldots7$ и
проверьте, что соседние отличаются одним битом. \stars{2} Постройте
преобразователь двоичного кода в код Грея из элементов XOR для трёх бит.
\stars{3} Соберите обратный преобразователь и убедитесь, что двойное
преобразование возвращает исходное число.

\textbf{В Logra.} \stars{1} Урок \emph{«Код Грея блоком»} (книга «Сравнение
и коды»). \stars{2} Вкладка \emph{Конструктор}: два блока XOR. Пример
\code{gray} (библиотека примеров): сравните \code{binary} и \code{gray} на
переходе $3\to4$. \stars{3} Пример \code{binary-onehot} как образец
преобразования «код $\to$ код».

\textbf{Ожидаемый результат.} При переходе $3\to4$ меняются три
бита двоичного кода и один бит кода Грея. Двойное преобразование
возвращает число.

\textbf{Как проверить.} Вкладка \emph{Диаграммы}: заведите вход $b$
счётчиком $0$--$7$ и сосчитайте, сколько разрядов Грея меняется на
каждом шаге; везде должно быть ровно одно изменение.
\end{diybox}

\section{Те же узлы на Verilog}
\label{sec:coders-verilog}

Все узлы этой главы удобно описывать не рисунком, а текстом (язык Verilog
подробно разобран в главе~\ref{ch:verilog}). Правила чтения короткие.
\code{assign} означает «выход всегда равен этому выражению».
Блок \code{always @*} с \code{case} или \code{if} читается как «таблица или
цепочка условий, по которой выбирается одно значение»; в нём пишем
блокирующее присваивание \code{=} и \emph{в каждой ветке присваиваем
каждому выходу значение}.

Как запустить любой пример этого раздела в Logra: вкладка \emph{Код},
вставить модуль (или открыть файл из \code{code/b4\_coders}), нажать
\emph{Синтезировать} (Ctrl+Enter); вкладка \emph{RTL} показывает полученную
схему, вкладка \emph{Диаграммы} --- перебор входов. Рядом с каждым модулем лежит тестбенч
\code{tb\_*.v}: он перебирает входы, сверяет выходы с эталоном и печатает
\code{OK}.

\subsection{Шифратор $4\to2$ через case}

Таблица шифратора превращается в \code{case} почти буквально: слева
унитарный код, справа номер линии.

\vfile{\codepath/b4_coders/enc4to2.v}{code/b4\_coders/enc4to2.v}

\textbf{Как читается.} \code{case (x)} сравнивает вход с четырьмя
образцами сверху вниз; первый совпавший задаёт \code{y}. Ветка
\code{default} получает всё остальное: ноль единиц и несколько единиц.
Без неё для этих кодов \code{y} осталось бы «как было», то есть появилась бы
память (см.\ подраздел об ошибках ниже). Сравним с версией из «вентилей»,
которую мы рисовали в начале раздела:

\vfile{\codepath/b4_coders/enc4to2_or.v}{code/b4\_coders/enc4to2\_or.v}

На унитарных кодах обе версии одинаковы. На $x=0110$ версия из ИЛИ даёт
$y=11$ (чужой номер, как мы и предсказывали), а версия через \code{case}
честно «не знает» и выдаёт $00$: поведение на запрещённых входах теперь
записано явно, а не получается случайно.

\vfile{\codepath/b4_coders/tb_enc4to2.v}{code/b4\_coders/tb\_enc4to2.v}

\textbf{Сценарий на вкладке «Диаграммы».} Задайте $x$ значениями
\code{0001, 0010, 0100, 1000}: $y$ должен пройти $0,1,2,3$ у обоих модулей;
затем $x=0110$ --- и найдите, чем различаются $y$.

\subsection{Приоритетный шифратор: if/else if и casez}

Приоритет в тексте --- это \emph{порядок проверок}. Первая проверка,
оказавшаяся истинной, побеждает, остальные уже не смотрятся.

\vfile{\codepath/b4_coders/prio4_if.v}{code/b4\_coders/prio4\_if.v}

\textbf{Как читается.} Сначала \code{v = 1} по умолчанию. Затем цепочка:
«если активен $x_3$, номер 3; иначе, если $x_2$, номер 2; \ldots; иначе
(ничего нет) номер 0 и $v=0$». Каждая ветка присваивает и \code{y}, и \code{v}:
иначе получилась бы защёлка. Та же функция через \code{casez}, где знак
\code{?} означает «здесь может быть что угодно», и это в точности
$\times$ из таблицы приоритетного шифратора:

\vfile{\codepath/b4_coders/prio4_casez.v}{code/b4\_coders/prio4\_casez.v}

Строки \code{casez} проверяются сверху вниз, как цепочка \code{if}: для
$x=0110$ первая строка не подходит (старший бит нулевой), вторая
\code{01??} подходит, и $y=2$. Две записи равнозначны; \code{casez} нагляднее
повторяет таблицу, \code{if} лучше, когда условия сложнее «маски».

\vfile{\codepath/b4_coders/tb_prio4.v}{code/b4\_coders/tb\_prio4.v}

\textbf{Сценарий на вкладке «Диаграммы».} Входы $x$ счётчиком $0$--$15$:
при $x=0$ выход $v=0$; при любой единице в $x_3$ всегда $y=3$; на
$x=0110$ получается $y=2$, и обе версии модуля дают одинаковые графики.

\subsection{Дешифраторы: assign и case}

Дешифратор $2\to4$ можно записать по одному \code{assign} на выход
(формулы из начала главы) или таблицей \code{case}:

\vfile{\codepath/b4_coders/dec2to4_assign.v}{code/b4\_coders/dec2to4\_assign.v}

\vfile{\codepath/b4_coders/dec2to4_case.v}{code/b4\_coders/dec2to4\_case.v}

\textbf{Как читается.} В первом модуле каждая строка --- произведение
$E\,\overline{a_1}\,\overline{a_0}$ и т.\,д., знак \code{\textasciitilde} --- НЕ.
Во втором сначала проверяется разрешение: при \code{en = 0} все выходы нули,
иначе \code{case (a)} выбирает строку таблицы. Рядом ---
дешифратор $3\to8$ одной строкой: единица сдвигается влево на \code{a}
позиций, а при \code{en = 0} выход нулевой; и его же сборка из двух блоков
$2\to4$ по рисунку~\ref{fig:dc38}.

\vfile{\codepath/b4_coders/dec3to8.v}{code/b4\_coders/dec3to8.v}

\vfile{\codepath/b4_coders/dec3to8_cascade.v}{code/b4\_coders/dec3to8\_cascade.v}

В каскадной версии младшие два бита номера идут в оба блока, а старший
бит разрешает один из них: \code{en \& \textasciitilde a[2]} включает нижнюю
половину ($y_0\ldots y_3$), \code{en \& a[2]} --- верхнюю. Это в точности
схема с рисунка~\ref{fig:dc38}.

\textbf{Сценарий.} Тестбенч \code{tb\_dec.v} сравнивает все четыре
описания с эталоном на всех 16 комбинациях $(E,a)$. На вкладке
\emph{Диаграммы} задайте $a$ счётчиком $0$--$7$ при $en=1$: единица
«бежит» по выходам \code{dec3to8}; у каскада на вкладке \emph{RTL} видно
два одинаковых блока и элементы \code{\&}, \code{\textasciitilde}.

\subsection{Семисегментный дешифратор}

Таблица~\ref{tab:7seg} превращается в \code{case} буквально строка в
строку; семь выходов собраны в одну шину \code{seg = \{g,f,e,d,c,b,a\}}
(младший бит --- сегмент $a$).

\vfile{\codepath/b4_coders/seg7.v}{code/b4\_coders/seg7.v}

\textbf{Как читается.} Для цифры~1 получается \code{7'b0000110}, то есть
горят $b$ и $c$, как в таблице. Ветка \code{default} закрывает коды
$10$--$15$ и гасит индикатор. Это то самое решение для безразличных состояний,
о котором мы говорили выше; вместо нулей можно было бы написать \code{7'bx},
и оптимизатор выбрал бы значения сам (см.\ главу~\ref{ch:minimize}).
Тестбенч \code{tb\_seg7.v} проверяет все 16 кодов.

\textbf{Сценарий.} Вкладка \emph{Диаграммы}: вход $x$ счётчиком $0$--$15$;
на цифрах $0$--$9$ сравните выход с таблицей, на $10$--$15$ он нулевой.
Вкладка \emph{RTL}: посмотрите, во что превратился этот \code{case}:
цифры не «разбираются» по семи отдельным схемам, это одна таблица.

\subsection{Преобразователи кодов: код Грея}

Прямое преобразование --- одна строка: каждый бит равен XOR двух соседних
двоичных бит. Обратное требует цепочки: каждый следующий бит зависит от уже
полученного.

\vfile{\codepath/b4_coders/bin2gray.v}{code/b4\_coders/bin2gray.v}

\vfile{\codepath/b4_coders/gray2bin.v}{code/b4\_coders/gray2bin.v}

\textbf{Как читается.} \code{b \textasciicircum\ (b >> 1)} сдвигает число на
бит вправо и складывает с исходным по XOR: $g_2=b_2$, $g_1=b_2\oplus b_1$,
$g_0=b_1\oplus b_0$. Во втором модуле выход \code{b[1]} зависит от
\code{b[2]}, то есть цепочка тянется от старшего бита. Тестбенч
\code{tb\_gray.v} проверяет, что двойное преобразование возвращает число
и что соседние коды (в том числе $7\to0$) различаются ровно одним битом.

\textbf{Сценарий.} Вкладка \emph{Диаграммы}: $b$ счётчиком $0$--$7$, выход
\code{g} на каждом шаге меняется ровно одним битом.

\subsection{Типичные ошибки на Verilog}

\textbf{Защёлка из-за неполного if.} Если в комбинационном блоке при
каких-то условиях выходу ничего не присвоено, схема «запоминает» прошлое
значение, то есть возникает защёлка (латч): комбинационная схема внезапно
получает память.

\begin{vcode}
// ОШИБКА: при a = 2 и a = 3 выходу y ничего не присвоено
module dec_latch (
    input      [1:0] a,
    output reg [3:0] y
);
    always @* begin
        if (a == 2'd0)      y = 4'b0001;
        else if (a == 2'd1) y = 4'b0010;
    end
endmodule
\end{vcode}

На вкладке \emph{RTL} вместо чистой комбинационной логики появится блок
защёлки, а в журнале --- предупреждение о ней. Лечение: присвоить значение
\emph{при любом условии}: добавить завершающий \code{else} или записать
значение по умолчанию в самом начале блока.

\begin{vcode}
// ИСПРАВЛЕНО: значение по умолчанию в начале блока
module dec_fixed (
    input      [1:0] a,
    output reg [3:0] y
);
    always @* begin
        y = 4'b0000;                       // по умолчанию все выходы нули
        if (a == 2'd0)      y = 4'b0001;
        else if (a == 2'd1) y = 4'b0010;
    end
endmodule
\end{vcode}

\textbf{Нет default в case.} Та же беда с \code{case}: если перечислены не
все значения селектора и нет \code{default}, для остальных значений выход
запоминается. Если все значения перечислены (как в $4$ строках для $2$-битного
\code{a}), \code{default} не обязателен, но писать его всё равно полезно.

\begin{vcode}
// ОШИБКА: для x, отличного от четырёх унитарных кодов, y запоминается
module enc_latch (
    input      [3:0] x,
    output reg [1:0] y
);
    always @* begin
        case (x)
            4'b0001: y = 2'd0;
            4'b0010: y = 2'd1;
            4'b0100: y = 2'd2;
            4'b1000: y = 2'd3;
        endcase
    end
endmodule
\end{vcode}

\textbf{Блокирующее и неблокирующее.} В комбинационном блоке
\code{always @*} используйте \code{=}. Оператор \code{<=} нужен для
триггеров (глава~\ref{ch:ff}); в комбинационной схеме он работает, но
путает читателя и ведёт к странным результатам, когда выход используется
в том же блоке.

\begin{diybox}[Допиши и измени модуль]
\textbf{Задание.} \stars{1} \emph{Допиши.} Опишите приоритетный шифратор
$8\to3$ с выходом $v$ через \code{casez}: достроить восемь строк по образцу
\code{prio4\_casez}. \stars{2} \emph{Измени.} Добавьте в
\code{seg7} шестнадцатеричные цифры $A\ldots F$ (строчные \code{b} и
\code{d}, чтобы не путать с $8$ и $0$). \stars{3} \emph{Этап проекта
«Охранная панель».} Напишите модуль \code{zone\_display}: четыре датчика зон
$z$ на входе, семь сегментов на выходе, при $z=0000$ индикатор погашен.

\textbf{В Logra.} Вкладка \emph{Код}: введите свой модуль, нажмите
\emph{Синтезировать}. Вкладка \emph{Диаграммы}: для \stars{1} задайте $x$
счётчиком $0$--$255$ (или несколько характерных значений: \code{00000000},
\code{00000001}, \code{00110000}, \code{10000000}); для \stars{2} --- $x$
счётчиком $0$--$15$; для \stars{3} --- шесть строк проверочной таблицы.

\textbf{Ожидаемый результат.} \stars{1} $x=00110000$ даёт $y=5$, $v=1$;
$x=0$ даёт $v=0$. \stars{2} Для $A$ получается \code{1110111}. \stars{3}
Для зон~1, 2, 3: \code{0000110}, \code{1011011}, \code{1001111}; для двух
зон показана старшая; при $z=0$ все нули.

\textbf{Как проверить.} Для заранее известных ответов сравните диаграмму с
ожиданиями выше; в папке \code{code/b4\_coders} лежит тестбенч
\code{tb\_samples.v}, проверяющий образцы ответов (он печатает \code{OK}).

\end{diybox}

\textbf{Образцы ответов к заданию «Допиши и измени модуль».}

\textbf{Образец ответа к \stars{1}.}
\vfile{\codepath/b4_coders/prio8_casez.v}{code/b4\_coders/prio8\_casez.v}

\textbf{Образец ответа к \stars{2}} (остальные строки таблицы те же, что в
\code{seg7}; здесь файл целиком):
\vfile{\codepath/b4_coders/seg7_hex.v}{code/b4\_coders/seg7\_hex.v}

\textbf{Образец ответа к \stars{3}.}
\vfile{\codepath/b4_coders/zone_display.v}{code/b4\_coders/zone\_display.v}

\section{Сквозной проект: номер зоны на индикаторе}

Соберём всё, что узнали, в одну схему охранной панели. Четыре датчика зон
$z_3z_2z_1z_0$ (зона~0 --- «ничего не сработало»). Приоритетный шифратор
превращает их в номер ($2$ бита). Индикатору нужна цифра в виде
четырёх бит, поэтому к номеру слева допишем два нуля. Дешифратор для
индикатора переводит цифру в семь сегментов.

\begin{figure}[H]
\centering
\begin{tikzpicture}[node distance=15mm, block/.append style={minimum width=22mm, font=\sffamily\footnotesize}]
  \node[block] (enc) {Приоритетный\\шифратор};
  \node[block, right=of enc] (cat) {Склейка\\\code{00}\,\&\,номер};
  \node[block, right=of cat] (seg) {BCD$\to$7\\сегментов};
  \node[block, right=of seg] (ind) {Индикатор\\$a\ldots g$};
  \draw[arr] ($(enc.west)+(-0.9,0)$) node[left] {$z$ (4)} -- (enc.west);
  \draw[arr] (enc.east) -- node[above, lbl] {2 бита} (cat.west);
  \draw[arr] (cat.east) -- node[above, lbl] {4 бита} (seg.west);
  \draw[arr] (seg.east) -- node[above, lbl] {7 бит} (ind.west);
\end{tikzpicture}
\caption{Этап~4а проекта «Охранная панель»: от четырёх датчиков зон к цифре на индикаторе}
\label{fig:zone_display}
\end{figure}

Что должно получиться (приоритет у старшей зоны):

\begin{table}[H]
\centering
\caption{Проверочная таблица номера зоны}
\begin{tabular}{c|c|c}
\toprule
Датчики $z_3z_2z_1z_0$ & Номер на индикаторе & Комментарий \\
\midrule
\code{0000} & 0 & спокойно \\
\code{0010} & 1 & сработала зона 1 \\
\code{0100} & 2 & сработала зона 2 \\
\code{1000} & 3 & сработала зона 3 \\
\code{0110} & 2 & две зоны: показываем старшую \\
\code{1010} & 3 & две зоны: показываем старшую \\
\bottomrule
\end{tabular}
\end{table}

\begin{diybox}[Охранная панель: этап 4а, номер зоны]
\textbf{Задание.} \stars{1} Соберите цепочку «приоритетный шифратор
$\to$ склейка $\to$ BCD$\to$7 сегментов» и проверьте все шесть
строк таблицы. \stars{2} Используйте выход $V$ шифратора, чтобы при
$z=0000$ индикатор гас, а не показывал «0» (каждый сегмент пропустите
через И с $V$). \stars{3} \emph{Этап проекта «Охранная панель»:} вместо
гашения покажите при спокойствии чёрточку «$-$» (горит только сегмент
$g$): сегменты $a\ldots f$ пропускайте через И с $V$, а сегмент $g$ ---
через ИЛИ с $\overline V$.

\textbf{В Logra.} Вкладка \emph{Конструктор}. Блоки «Приоритетный
шифратор» (ширина 4), «Объединение шин», «Константа» и «BCD\,$\to$\,7
сегментов» из группы «Выбор». Образцы: уроки \emph{«Приоритетный
шифратор»} и \emph{«Цифра на семи сегментах»}; пример
\code{segment-decoder} (библиотека примеров).

\textbf{Ожидаемый результат.} На выходе из семи бит: для зоны~1 код
\code{0000110}, для зоны~2 --- \code{1011011}, для зоны~3 ---
\code{1001111}; для зоны~0 --- \code{0111111} (в варианте \stars{2}
все нули, в варианте \stars{3} \code{1000000}).

\textbf{Как проверить.} Критерий приёмки: все шесть строк таблицы дают
ожидаемые коды на вкладке \emph{Диаграммы}. На вкладке \emph{RTL} видно
три основных блока. Результат этого этапа пригодится в главе, где
мы научим схему сравнивать числа.
\end{diybox}

\begin{warnbox}[Типичные ошибки]
\begin{itemize}
  \item \textbf{Перепутать порядок разрядов.} В дешифраторе $a_0$ --- младший
        бит номера. Если поменять $a_0$ и $a_1$, зажжётся не та линия, но
        схема будет «исправной». Проверьте на номере 1 ($01$), а не на 0 или 3.
  \item \textbf{Забыть вход разрешения.} При $E=0$ все выходы равны нулю.
        Если схема «молчит», первым делом проверьте $E$.
  \item \textbf{Простой шифратор там, где входов может быть несколько.}
        Две единицы на входе дают чужой номер. Если одновременные события
        возможны, нужен приоритетный шифратор, и $V$ для случая «ничего нет».
  \item \textbf{Индикатор показывает «негатив».} Общий анод: сегмент
        горит при нуле, а вы подаёте единицу. Нужны инверторы на семи
        выходах.
  \item \textbf{Неполный \code{if} или \code{case} в Verilog.} Если при
        каком-то входе выходу ничего не присвоено, получается защёлка.
        Присваивайте значение по умолчанию в начале блока \code{always @*} или
        пишите \code{else} и \code{default}.
  \item \textbf{\code{<=} в комбинационном блоке.} Для комбинационной
        логики пишите \code{=}; \code{<=} оставьте для триггеров.
\end{itemize}
\end{warnbox}

\begin{ideabox}
\textbf{Дешифратор}: номер $\to$ единица на одном из $2^n$ выходов.
\textbf{Шифратор}: единица на одном из входов $\to$ её номер.
\textbf{Приоритетный шифратор} при нескольких активных входах выбирает старший.
\textbf{Код Грея} меняет один бит на шаг. Цифровой индикатор --- тоже
дешифратор, только выход у него не «одна горячая линия», а набор сегментов.
\end{ideabox}

\section{Вопросы и задачи}
\label{sec:tasks-coders}

Ответы собраны в приложении~\ref{app:answers}.

\begin{taskbox}[Задачи к главе]
\begin{enumerate}
  \item \stars{1}~Дешифратор $3\to 8$: сколько у него входов и выходов? Какая линия активна, если на входе 101?
  \item \stars{1}~Запишите выражения для всех четырёх выходов дешифратора $2\to 4$ с входом разрешения EN.
  \item \stars{1}~Какие сегменты семисегментного индикатора должны гореть для цифры 7 (сегменты $a$--$g$ в стандартной раскладке)?
  \item \stars{1}~У приоритетного шифратора $4\to 2$ входы $I_3I_2I_1I_0 = 0110$. Какой код на выходе и почему?
  \item \stars{1}~Запишите число 6 в двоичном коде и в коде Грея. Чем код Грея удобен для датчиков угла?
  \item \stars{2}~Сколько дешифраторов $2\to 4$ понадобится, чтобы собрать дешифратор $4\to 16$?
  \item \stars{2}~Найдите минимальную форму сегмента $c$ семисегментного индикатора, считая коды $10$--$15$ безразличными. Сколько элементов потребуется?
  \item \stars{2}~Переведите код Грея \code{110} в двоичный код. Затем проверьте результат обратным преобразованием.
  \item \stars{2}~\emph{Найдите ошибку.} Ученик собрал шифратор $4\to2$ из двух элементов ИЛИ ($y_0=x_1+x_3$, $y_1=x_2+x_3$) и получил для входов $x_3x_2x_1x_0=0110$ код $11$. Он считает, что должно быть $10$. Кто прав? Что нужно изменить в схеме?
  \item \stars{2}~\emph{Найдите ошибку.} Для числа 5 ученик записал код Грея \code{100}, «просто сдвинув число». Чему равен настоящий код Грея числа 5? Какому числу соответствует \code{100}?
  \item \stars{3}~\emph{Проект с критерием (этап проекта «Охранная панель», 4а).} Постройте в Конструкторе цепочку «приоритетный шифратор $\to$ склейка $\to$ BCD$\to$7 сегментов» для четырёх датчиков зон. Критерий приёмки: все строки проверочной таблицы (нулевая зона, зоны 1, 2, 3, две зоны одновременно) дают верный код сегментов на вкладке Диаграммы.
  \item \stars{3}~\emph{Проект с критерием.} Соберите дешифратор $3\to8$ из двух $2\to4$ с входом разрешения. Критерий приёмки: при каждом из восьми номеров активна ровно одна линия, при $E=0$ все выходы нулевые.
  \item \stars{3}~\emph{Проект с критерием.} Соберите преобразователь $3$-битного двоичного кода в код Грея и обратно. Критерий приёмки: при двойном преобразовании выход равен входу на всех восьми значениях, а на выходе прямого преобразователя при счёте $0$--$7$ на каждом шаге меняется ровно один бит.
  \item \stars{2}~\emph{Найдите ошибку (Verilog).} Ученик описал приоритетный шифратор цепочкой \code{if (x[3]) y = 2'd3; else if (x[2]) y = 2'd2; else if (x[1]) y = 2'd1; else if (x[0]) y = 2'd0;} внутри \code{always @*}, без последнего \code{else}. После синтеза на вкладке RTL появился блок защёлки. Почему и как исправить?
  \item \stars{3}~\emph{Проект с критерием.} Опишите дешифратор $3\to8$ с входом разрешения на Verilog двумя способами: одной строкой со сдвигом и из двух дешифраторов $2\to4$ (как на рисунке~\ref{fig:dc38}). Критерий приёмки: оба модуля совпадают на всех 16 комбинациях $(en,a)$ на вкладке Диаграммы, при $en=0$ выходы нулевые.
\end{enumerate}
\end{taskbox}

\section{Что дальше}

Теперь мы умеем превращать номер в «одну горячую линию» и обратно, показывать
цифру на индикаторе и безопасно считывать положение датчика. Осталась
родственная задача: не просто зажечь одну линию, а \emph{выбрать} один из
нескольких сигналов и пропустить его на единственный выход. Сегодняшний
дешифратор, скомбинированный с элементами И, как раз превращается в такой
выбор --- мультиплексор. Ему посвящена следующая глава.
